% =====================================================================
%  tikzlib · matematicas/recta-numerica
%  Requiere: \usepackage{tikz} ; \usetikzlibrary{arrows.meta}
%  \rectanumerica{min}{max}        → recta con marcas enteras
%  \rectapunto{min}{max}{p}        → igual, con un punto marcado en p
%  Preview: tikzlib/previews/recta-numerica.png
% =====================================================================
\providecolor{accent}{HTML}{1A8917}

\newcommand{\rectanumerica}[2]{%
\begin{tikzpicture}[x=1cm, font=\sffamily\small]
  \draw[-{Stealth[length=4pt]}, line width=0.9pt] (#1-0.6,0) -- (#2+0.7,0);
  \foreach \n in {#1,...,#2}{%
    \draw[line width=0.7pt] (\n,0.13) -- (\n,-0.13);
    \node[below=1pt] at (\n,-0.13) {\n};}
\end{tikzpicture}}

\newcommand{\rectapunto}[3]{%
\begin{tikzpicture}[x=1cm, font=\sffamily\small]
  \draw[-{Stealth[length=4pt]}, line width=0.9pt] (#1-0.6,0) -- (#2+0.7,0);
  \foreach \n in {#1,...,#2}{%
    \draw[line width=0.7pt] (\n,0.13) -- (\n,-0.13);
    \node[below=1pt] at (\n,-0.13) {\n};}
  \fill[accent] (#3,0) circle (2.6pt);
\end{tikzpicture}}
